\chapter{Nota puente entre admisión, demo y claim}

\section{Función del puente}
\begin{observation}
La admisión metodológica, la demo seleccionada y la futura traducción a claim no deben leerse como piezas aisladas. Forman una secuencia disciplinada dentro del Tomo III.
\end{observation}

\begin{observation}
La admisión metodológica fija la puerta de entrada. La demo seleccionada fija la pieza curada que merece exposición y seguimiento. La traducción a claim fija la posibilidad de convertir esa pieza en una formulación más estricta, verificable y comunicable.
\end{observation}

\section{Secuencia mínima}
\begin{definition}
La secuencia mínima de tratamiento en Tomo III puede leerse así:
\[
\text{admisión metodológica}
\;\longrightarrow\;
\text{demo seleccionada}
\;\longrightarrow\;
\text{traducción disciplinada a claim}.
\]
\end{definition}

\begin{observation}
Cada tramo presupone el anterior. No toda pieza admitida debe convertirse en demo, y no toda demo debe traducirse de inmediato a claim.
\end{observation}

\section{Condición de continuidad}
\begin{definition}
Existe continuidad metodológica suficiente cuando una pieza conserva, a lo largo de esta secuencia, dependencia canónica visible, trazabilidad legible y función reconocible para el propósito declarado.
\end{definition}

\begin{observation}
Si alguno de estos elementos se rompe, la secuencia debe detenerse o devolver la pieza a una fase previa de curación.
\end{observation}

\section{Condición mínima de interrupción}
\begin{definition}
La secuencia se interrumpe legítimamente cuando:
\begin{enumerate}
  \item la pieza admitida no alcanza claridad suficiente para funcionar como demo;
  \item la demo seleccionada no sostiene una traducción disciplinada sin pérdida excesiva;
  \item el intento de claim exige supuestos que todavía no han sido fijados en el canon.
\end{enumerate}
\end{definition}

\begin{observation}
La interrupción no debe leerse como fracaso del sistema, sino como protección metodológica frente a inflación prematura.
\end{observation}

\section{Cláusula de alcance}
\begin{observation}
Esta nota puente no define todavía el formato final de claim ni su auditoría exhaustiva. Su función es fijar la continuidad mínima entre entrada metodológica, selección curatorial y proyección disciplinada.
\end{observation}

